\paragraph{Flat coefficients over a Noetherian base, in every characteristic.}

\begin{proposition}\label{prop:DE-E2-flat-Noetherian}
\status{Proved; root and monodromy full independent proof/source audits passed.}
Let $S$ be locally Noetherian and let $X\to S$ be flat and locally of finite presentation.  Restriction is an equivalence between quasi-coherent crystals whose values on infinitesimal tests are flat modules and ordinary stratified quasi-coherent $\OO_X$-modules that are flat over $\OO_X$.  There is no finite-generation hypothesis on the coefficients, and no restriction on the characteristic of $S$.

This equivalence restricts to finite locally free coefficients.  It holds for either the small or big infinitesimal site, with arbitrary classical test schemes.
\end{proposition}
\begin{proof}
We state explicitly the relative derived comparison being used.  Zariski locally, $S$ and $X$ are affine Noetherian schemes and $X/S$ is separated, of finite type and of finite Tor-amplitude.  Jiang's Theorem~2.1 and Remark~2.4 identify the category of derived stratifications on the intrinsic formal diagonal nerve with modules over his derived differential-operator algebra.  His Theorem~6.3 identifies the latter category with derived quasi-coherent crystals on $(X/S)_{\rm inf}$, compatibly with evaluation on $X$; Appendix~B, Theorem~B.5, identifies the small and big versions \citep[pp.~16--17, 32--34 and 43--45]{Jiang2023Derived}.

For clarity, the formal-diagonal interpretation uses Jiang's Proposition~6.2: on an affine closed immersion with ideal $J$, the category of derived sheaves on the formal completion is
\[
       \varprojlim_n D(B/J^n),
\]
with derived pullback transitions.  Equivalently it is the category of derived $J$-complete objects, or, by the local-cohomology/completion equivalence, the category denoted $\Gamma_J(\mathrm{QCoh}(\Spec B))$ in his Theorem~2.1.  The pullbacks in that theorem are the corresponding supported pullbacks.  Remark~2.4 gives precisely the resulting derived stratifying category.  We use this comparison only after checking that the ordinary data define objects and morphisms in these derived categories.  Flatness of $X/S$ ensures that its iterated products over $S$ here are the ordinary products, without extra derived structure.

Let $M$ be an ordinary flat stratified $\OO_X$-module.  On every finite diagonal neighbourhood, all its pullbacks are flat over that neighbourhood.  Their ordinary pullbacks therefore equal their derived pullbacks.  The given compatible isomorphisms and cocycle conditions define an object on the derived formal diagonal nerve.  Mapping spaces between the ordinary modules at each level have no positive homotopy groups, so these ordinary cocycle identities supply the required coherent identifications.  The cited comparison supplies a derived crystal $\mathcal K$ with $\mathcal K_X=M[0]$.

We now prove that its values are ordinary flat modules; this is a separate step from the cited comparison.  Let $B$ be any ring, let $I^N=0$, and let $K\in D(B)$ satisfy
\[
          K\otimes_B^{\mathbf L}B/I\simeq P[0],
          \qquad P\text{ flat over }B/I.
\]
For every $B/I$-module $Q$, tensor associativity gives
$K\otimes_B^{\mathbf L}Q\simeq(P\otimes_{B/I}Q)[0]$.
Filter an arbitrary $B$-module $Q'$ by the finite filtration $I^jQ'$.  Each successive quotient is a $B/I$-module, so the associated exact triangles show that $K\otimes_B^{\mathbf L}Q'$ is concentrated in degree zero.  Taking $Q'=B$ proves that $K$ is an ordinary module, and varying $Q'$ proves that it is flat.  This argument allows unbounded $K$ and imposes no Noetherian condition on $B$.

For a test $U\hookrightarrow T$, apply this argument to $K=\mathcal K_T$, whose derived reduction is the flat module pulled back from $M$ on $U$.  Thus $\mathcal K$ is an ordinary flat crystal.  Conversely an ordinary flat crystal has derived Cartesian transition maps, so defines a derived crystal.  The two constructions are inverse by the cited comparison, and the same reasoning on the zero-truncated mapping spaces identifies their ordinary morphisms.

These arguments initially use Jiang's finite-type classical tests.  They extend to arbitrary tests without an assumption about infinite completed tensor products.  Indeed, locally a test $\Spec C\hookrightarrow\Spec B$ mapping to an affine $\Spec A\subset X$, with $A=R[z_1,\ldots,z_r]/K$, maps to a finite-type test
\[
 \Spec A\hookrightarrow\Spec\bigl(R[z_1,\ldots,z_r]/K^N\bigr)
\]
by lifting the $z_i$, for sufficiently large $N$.  Two such choices have a common refinement: use the finite-type product of their ambient thickenings, then a sufficiently high infinitesimal neighbourhood of the diagonal copy of $\Spec A$.  Parallel compatible maps are equalized by their closed equalizer, which still contains that copy of $\Spec A$.  Pullback of the finite-type crystal is therefore independent of all choices and gives the unique extension, on objects and morphisms.  This also proves compatibility with affine localization and gluing.

Finally, if $M$ is finite locally free, nilpotent Nakayama shows that every test value is finitely generated.  At a local test ring, lift a basis of its reduction and use flatness to identify the kernel modulo the nilpotent ideal; a second nilpotent Nakayama argument makes that kernel zero.  The test value is consequently finite locally free.  The equivalence thus restricts as asserted.
\end{proof}

\begin{corollary}\label{cor:DE-E2-coherent-any-field}
\status{Proved; root and monodromy full independent proof/source audits passed.}
For every field $k$ and every $k$-scheme $X$ locally of finite type, restriction gives an equivalence between coherent infinitesimal crystals and ordinary coherent stratified $\OO_X$-modules.  Smoothness, reducedness and $F$-purity are unnecessary.  Every coefficient module in either category is locally free of finite rank.
\end{corollary}
\begin{proof}
A coherent stratified module is locally free.  After the faithfully flat extension to an algebraic closure, fix a closed point $x$ and pull its stratification on each diagonal neighbourhood back along the second factor $x$.  The compatible isomorphisms identify the completed stalk with
\[
        \widehat{\OO}_{X,x}\otimes_k M_x/\mathfrak m_xM_x.
\]
Faithful flatness of the Noetherian local completion then makes the stalk free.  The nonfree locus of a finitely presented module is closed, so checking closed points on each affine finite-type open suffices.  Local freeness descends along the field extension.  This is also the argument behind Tonini--Zhang's Theorem~6.8(3).

For coherent infinitesimal crystals, Tonini--Zhang's Theorem~6.8(4), with its proof using Lemmas~6.10 and~6.15, shows that every test value is finite locally free \citep{ToniniZhang2019}.  The statement includes arbitrary ground fields.  The proposition now applies, since a scheme locally of finite type over a field is flat over that field, and proves the asserted equivalence.
\end{proof}

\begin{remark}[Scope and the subsequent coherent comparison]\label{rem:DE-E2-flat-Noetherian-scope}
\status{Proved scope statement.}
The proposition does not replace ordinary nonflat coherent coefficients by flat ones over a general base.  Such coefficients occur even for smooth morphisms: the pullback of a coherent module from the base has its constant-base crystal, regardless of its flatness.  Theorem~\ref{thm:DE-E2-coherent-Noetherian} below separately proves fullness and effectivity for nonflat coherent coefficients over every locally Noetherian base.  The present proposition has a different strength: it allows arbitrary, possibly infinitely generated, flat coefficients.  The original arbitrary-module assertion in Part~E.2 remains refuted by the explicit examples above.
\end{remark}
